\chapter{Set-Theoretic Foundations}

\section{Algebras of Sets}

\begin{definition}[Algebra of Sets]
Let $X$ be a set. A collection $\mathcal{A} \subseteq \Powerset{X}$ is an \textbf{algebra} (or \textbf{field}) of sets if:
\begin{enumerate}
    \item $\emptyset \in \mathcal{A}$
    \item If $A \in \mathcal{A}$, then $A^c = X \setminus A \in \mathcal{A}$
    \item If $A, B \in \mathcal{A}$, then $A \cup B \in \mathcal{A}$
\end{enumerate}
\end{definition}

\begin{proposition}[Properties of Algebras]
Let $\mathcal{A}$ be an algebra on $X$. Then:
\begin{enumerate}
    \item $X \in \mathcal{A}$
    \item If $A, B \in \mathcal{A}$, then $A \cap B \in \mathcal{A}$
    \item If $A, B \in \mathcal{A}$, then $A \setminus B \in \mathcal{A}$
    \item For any finite collection $A_1, \ldots, A_n \in \mathcal{A}$, we have $\bigcup_{i=1}^n A_i \in \mathcal{A}$ and $\bigcap_{i=1}^n A_i \in \mathcal{A}$
\end{enumerate}
\end{proposition}

\begin{proof}
(1) $X = \emptyset^c \in \mathcal{A}$.

(2) $A \cap B = (A^c \cup B^c)^c \in \mathcal{A}$ by De Morgan.

(3) $A \setminus B = A \cap B^c \in \mathcal{A}$.

(4) Induction on $n$ using axiom 3 and property 2.
\end{proof}

\section{$\sigma$-Algebras}

\begin{definition}[$\sigma$-Algebra]
Let $X$ be a set. A collection $\mathcal{M} \subseteq \Powerset{X}$ is a \textbf{$\sigma$-algebra} if:
\begin{enumerate}
    \item $\emptyset \in \mathcal{M}$
    \item If $A \in \mathcal{M}$, then $A^c \in \mathcal{M}$
    \item If $(A_n)_{n \in \bb{N}} \subseteq \mathcal{M}$, then $\bigcup_{n=1}^{\infty} A_n \in \mathcal{M}$
\end{enumerate}
\end{definition}

\begin{remark}
Every $\sigma$-algebra is an algebra, but not conversely.
\end{remark}

\begin{proposition}[Properties of $\sigma$-Algebras]
Let $\mathcal{M}$ be a $\sigma$-algebra on $X$. Then:
\begin{enumerate}
    \item $X \in \mathcal{M}$
    \item If $(A_n)_{n \in \bb{N}} \subseteq \mathcal{M}$, then $\bigcap_{n=1}^{\infty} A_n \in \mathcal{M}$
    \item If $A, B \in \mathcal{M}$, then $A \setminus B \in \mathcal{M}$
\end{enumerate}
\end{proposition}

\begin{proof}
(1) $X = \emptyset^c \in \mathcal{M}$.

(2) $\bigcap_{n=1}^{\infty} A_n = \left(\bigcup_{n=1}^{\infty} A_n^c\right)^c \in \mathcal{M}$ by De Morgan.

(3) $A \setminus B = A \cap B^c \in \mathcal{M}$.
\end{proof}

\begin{definition}[Generated $\sigma$-Algebra]
Let $\mathcal{F} \subseteq \Powerset{X}$. The \textbf{$\sigma$-algebra generated by $\mathcal{F}$}, denoted $\GenSigmaAlg{\mathcal{F}}$, is:
\[ \GenSigmaAlg{\mathcal{F}} = \bigcap \{ \mathcal{M} \mid \mathcal{M} \text{ is a } \sigma\text{-algebra on } X \text{ and } \mathcal{F} \subseteq \mathcal{M} \} \]
\end{definition}

\begin{proposition}
$\GenSigmaAlg{\mathcal{F}}$ is the smallest $\sigma$-algebra containing $\mathcal{F}$.
\end{proposition}

\begin{proof}
The intersection of $\sigma$-algebras is a $\sigma$-algebra. Since $\Powerset{X}$ is a $\sigma$-algebra containing $\mathcal{F}$, the intersection is nonempty. By construction, $\GenSigmaAlg{\mathcal{F}}$ contains $\mathcal{F}$ and is contained in every $\sigma$-algebra containing $\mathcal{F}$.
\end{proof}

\begin{definition}[Borel $\sigma$-Algebra]
Let $(X, \mathcal{T})$ be a topological space. The \textbf{Borel $\sigma$-algebra} is $\BorelAlg{X} = \GenSigmaAlg{\mathcal{T}}$, i.e., the $\sigma$-algebra generated by the open sets. Elements of $\BorelAlg{X}$ are called \textbf{Borel sets}.
\end{definition}

\begin{proposition}
On $\bb{R}$, the Borel $\sigma$-algebra can be generated by:
\begin{enumerate}
    \item Open intervals: $\GenSigmaAlg{\{(a, b) \mid a, b \in \bb{R}\}} = \BorelAlg{\bb{R}}$
    \item Half-open intervals: $\GenSigmaAlg{\{(a, b] \mid a, b \in \bb{R}\}} = \BorelAlg{\bb{R}}$
    \item Rays: $\GenSigmaAlg{\{(a, \infty) \mid a \in \bb{R}\}} = \BorelAlg{\bb{R}}$
\end{enumerate}
\end{proposition}

\begin{proof}
(1) Every open set in $\bb{R}$ is a countable union of open intervals. Thus $\GenSigmaAlg{\text{open intervals}} \supseteq \GenSigmaAlg{\text{open sets}} = \BorelAlg{\bb{R}}$. The reverse is clear since open intervals are open.

(2) $(a, b] = \bigcap_{n=1}^\infty (a, b + 1/n)$, so half-open intervals are Borel. Conversely, $(a, b) = \bigcup_{n=1}^\infty (a, b - 1/n]$.

(3) $(a, \infty)$ is open, so rays are Borel. Conversely, $(a, b) = (a, \infty) \setminus [b, \infty) = (a, \infty) \cap (-\infty, b)$ where $(-\infty, b) = \bigcup_{n=1}^\infty (-n, b)$ and $(-n, b) = (-n, \infty) \setminus [b, \infty)$.
\end{proof}

\section{$\pi$-Systems}

\begin{definition}[$\pi$-System]
A collection $\mathcal{P} \subseteq \Powerset{X}$ is a \textbf{$\pi$-system} if it is closed under finite intersections:
\[ A, B \in \mathcal{P} \implies A \cap B \in \mathcal{P} \]
\end{definition}

\begin{example}
The collection of intervals $(a, b]$ on $\bb{R}$ forms a $\pi$-system.
\end{example}

\section{$\lambda$-Systems (Dynkin Systems)}

\begin{definition}[$\lambda$-System]
A collection $\mathcal{L} \subseteq \Powerset{X}$ is a \textbf{$\lambda$-system} (or \textbf{Dynkin system}) if:
\begin{enumerate}
    \item $X \in \mathcal{L}$
    \item If $A, B \in \mathcal{L}$ and $A \subseteq B$, then $B \setminus A \in \mathcal{L}$
    \item If $(A_n)_{n \in \bb{N}} \subseteq \mathcal{L}$ with $A_n \subseteq A_{n+1}$ for all $n$, then $\bigcup_{n=1}^{\infty} A_n \in \mathcal{L}$
\end{enumerate}
\end{definition}

\begin{proposition}
Every $\sigma$-algebra is a $\lambda$-system.
\end{proposition}

\begin{proof}
Let $\mathcal{M}$ be a $\sigma$-algebra.

(1) $X \in \mathcal{M}$.

(2) If $A \subseteq B$ with $A, B \in \mathcal{M}$, then $B \setminus A = B \cap A^c \in \mathcal{M}$.

(3) If $A_n \nearrow$, then $\bigcup A_n \in \mathcal{M}$ since $\mathcal{M}$ is closed under all countable unions.
\end{proof}

\begin{definition}[Generated $\lambda$-System]
For $\mathcal{F} \subseteq \Powerset{X}$, let $\GenLambda{\mathcal{F}}$ denote the smallest $\lambda$-system containing $\mathcal{F}$.
\end{definition}

\section{Monotone Classes}

\begin{definition}[Monotone Class]
A collection $\mathcal{N} \subseteq \Powerset{X}$ is a \textbf{monotone class} if:
\begin{enumerate}
    \item If $(A_n)_{n \in \bb{N}} \subseteq \mathcal{N}$ with $A_n \subseteq A_{n+1}$, then $\bigcup_{n=1}^{\infty} A_n \in \mathcal{N}$
    \item If $(A_n)_{n \in \bb{N}} \subseteq \mathcal{N}$ with $A_n \supseteq A_{n+1}$, then $\bigcap_{n=1}^{\infty} A_n \in \mathcal{N}$
\end{enumerate}
\end{definition}

\begin{proposition}
Every $\sigma$-algebra is a monotone class.
\end{proposition}

\begin{proof}
Let $\mathcal{M}$ be a $\sigma$-algebra. If $A_n \nearrow A$, then $A = \bigcup A_n \in \mathcal{M}$. If $A_n \searrow A$, then $A = \bigcap A_n \in \mathcal{M}$.
\end{proof}

\begin{definition}[Generated Monotone Class]
For $\mathcal{F} \subseteq \Powerset{X}$, let $\GenMonotone{\mathcal{F}}$ denote the smallest monotone class containing $\mathcal{F}$.
\end{definition}

\section{Dynkin's $\pi$-$\lambda$ Theorem}

\begin{theorem}[Dynkin's $\pi$-$\lambda$ Theorem]
If $\mathcal{P}$ is a $\pi$-system and $\mathcal{L}$ is a $\lambda$-system with $\mathcal{P} \subseteq \mathcal{L}$, then $\GenSigmaAlg{\mathcal{P}} \subseteq \mathcal{L}$.

Equivalently: If $\mathcal{P}$ is a $\pi$-system, then $\GenLambda{\mathcal{P}} = \GenSigmaAlg{\mathcal{P}}$.
\end{theorem}

\begin{proof}
We show that $\GenLambda{\mathcal{P}}$ is a $\sigma$-algebra. Since every $\sigma$-algebra is a $\lambda$-system, we have $\GenSigmaAlg{\mathcal{P}} \subseteq \GenLambda{\mathcal{P}}$.

For the reverse inclusion, it suffices to show $\GenLambda{\mathcal{P}}$ is closed under finite intersections (making it both a $\pi$-system and $\lambda$-system, hence a $\sigma$-algebra).

Define $\mathcal{G}_1 = \{A \in \GenLambda{\mathcal{P}} \mid A \cap P \in \GenLambda{\mathcal{P}} \text{ for all } P \in \mathcal{P}\}$.

\textbf{Claim 1:} $\mathcal{G}_1$ is a $\lambda$-system.
\begin{itemize}
    \item $X \in \mathcal{G}_1$: For any $P \in \mathcal{P}$, $X \cap P = P \in \mathcal{P} \subseteq \GenLambda{\mathcal{P}}$.
    \item If $A \subseteq B$ with $A, B \in \mathcal{G}_1$: For any $P \in \mathcal{P}$, $(B \setminus A) \cap P = (B \cap P) \setminus (A \cap P) \in \GenLambda{\mathcal{P}}$ since $\GenLambda{\mathcal{P}}$ is a $\lambda$-system and $A \cap P \subseteq B \cap P$.
    \item If $A_n \nearrow A$ with $A_n \in \mathcal{G}_1$: For any $P$, $A_n \cap P \nearrow A \cap P$, so $A \cap P \in \GenLambda{\mathcal{P}}$.
\end{itemize}

Since $\mathcal{P}$ is a $\pi$-system, $\mathcal{P} \subseteq \mathcal{G}_1$. Since $\mathcal{G}_1$ is a $\lambda$-system containing $\mathcal{P}$, we have $\GenLambda{\mathcal{P}} \subseteq \mathcal{G}_1$, so $\GenLambda{\mathcal{P}} = \mathcal{G}_1$.

Define $\mathcal{G}_2 = \{A \in \GenLambda{\mathcal{P}} \mid A \cap B \in \GenLambda{\mathcal{P}} \text{ for all } B \in \GenLambda{\mathcal{P}}\}$.

\textbf{Claim 2:} $\mathcal{G}_2$ is a $\lambda$-system (proof similar to Claim 1).

Since $\GenLambda{\mathcal{P}} = \mathcal{G}_1$, for any $P \in \mathcal{P}$ and any $B \in \GenLambda{\mathcal{P}}$, we have $P \cap B \in \GenLambda{\mathcal{P}}$. Thus $\mathcal{P} \subseteq \mathcal{G}_2$, so $\GenLambda{\mathcal{P}} = \mathcal{G}_2$.

This means: for all $A, B \in \GenLambda{\mathcal{P}}$, $A \cap B \in \GenLambda{\mathcal{P}}$. So $\GenLambda{\mathcal{P}}$ is a $\pi$-system.

A collection that is both a $\pi$-system and $\lambda$-system is a $\sigma$-algebra.
\end{proof}

\section{Monotone Class Lemma}

\begin{theorem}[Monotone Class Lemma]
If $\mathcal{A}$ is an algebra on $X$, then $\GenSigmaAlg{\mathcal{A}} = \GenMonotone{\mathcal{A}}$.
\end{theorem}

\begin{proof}
Since $\GenSigmaAlg{\mathcal{A}}$ is a monotone class, $\GenMonotone{\mathcal{A}} \subseteq \GenSigmaAlg{\mathcal{A}}$.

For the reverse, we show $\GenMonotone{\mathcal{A}}$ is a $\sigma$-algebra.

\textbf{Step 1:} $\GenMonotone{\mathcal{A}}$ is closed under complements.

Define $\mathcal{G} = \{A \in \GenMonotone{\mathcal{A}} \mid A^c \in \GenMonotone{\mathcal{A}}\}$.

We verify $\mathcal{G}$ is a monotone class:
\begin{itemize}
    \item If $A_n \nearrow A$ with $A_n \in \mathcal{G}$: Then $A_n^c \searrow A^c$ with $A_n^c \in \GenMonotone{\mathcal{A}}$, so $A^c \in \GenMonotone{\mathcal{A}}$.
    \item If $A_n \searrow A$ with $A_n \in \mathcal{G}$: Then $A_n^c \nearrow A^c$, so $A^c \in \GenMonotone{\mathcal{A}}$.
\end{itemize}

Since $\mathcal{A} \subseteq \mathcal{G}$ and $\mathcal{G}$ is a monotone class, $\GenMonotone{\mathcal{A}} \subseteq \mathcal{G}$, so $\GenMonotone{\mathcal{A}} = \mathcal{G}$.

\textbf{Step 2:} $\GenMonotone{\mathcal{A}}$ is closed under finite unions.

For fixed $B \in \mathcal{A}$, define $\mathcal{H}_B = \{A \in \GenMonotone{\mathcal{A}} \mid A \cup B \in \GenMonotone{\mathcal{A}}\}$.

$\mathcal{H}_B$ is a monotone class (if $A_n \nearrow A$ then $A_n \cup B \nearrow A \cup B$). Since $\mathcal{A} \subseteq \mathcal{H}_B$, we get $\GenMonotone{\mathcal{A}} = \mathcal{H}_B$.

Now for fixed $A \in \GenMonotone{\mathcal{A}}$, define $\mathcal{H}_A = \{B \in \GenMonotone{\mathcal{A}} \mid A \cup B \in \GenMonotone{\mathcal{A}}\}$.

By the above, $\mathcal{A} \subseteq \mathcal{H}_A$. Since $\mathcal{H}_A$ is a monotone class, $\GenMonotone{\mathcal{A}} = \mathcal{H}_A$.

Thus $\GenMonotone{\mathcal{A}}$ is an algebra.

\textbf{Step 3:} An algebra that is also a monotone class is a $\sigma$-algebra.

Given $(A_n)_n \subseteq \GenMonotone{\mathcal{A}}$, define $B_n = \bigcup_{i=1}^n A_i \in \GenMonotone{\mathcal{A}}$. Then $B_n \nearrow \bigcup A_n$, so $\bigcup A_n \in \GenMonotone{\mathcal{A}}$.
\end{proof}


\chapter{Contents}

\section{Definition of Content}

\begin{definition}[Content]
Let $\mathcal{A}$ be an algebra on $X$. A function $\mu: \mathcal{A} \to [0, \infty]$ is a \textbf{content} (or \textbf{finitely additive measure}) if:
\begin{enumerate}
    \item $\mu(\emptyset) = 0$
    \item \textbf{Finite additivity:} If $A, B \in \mathcal{A}$ are disjoint, then $\mu(A \cup B) = \mu(A) + \mu(B)$
\end{enumerate}
\end{definition}

\begin{proposition}[Properties of Contents]
Let $\mu$ be a content on an algebra $\mathcal{A}$. Then:
\begin{enumerate}
    \item \textbf{Monotonicity:} $A \subseteq B \implies \mu(A) \leq \mu(B)$
    \item \textbf{Subtractivity:} If $A \subseteq B$ and $\mu(A) < \infty$, then $\mu(B \setminus A) = \mu(B) - \mu(A)$
    \item \textbf{Finite subadditivity:} $\mu(A \cup B) \leq \mu(A) + \mu(B)$
\end{enumerate}
\end{proposition}

\begin{proof}
(1) $B = A \cup (B \setminus A)$ disjoint, so $\mu(B) = \mu(A) + \mu(B \setminus A) \geq \mu(A)$.

(2) $\mu(B) = \mu(A) + \mu(B \setminus A)$, rearrange.

(3) $A \cup B = A \cup (B \setminus A)$, so $\mu(A \cup B) = \mu(A) + \mu(B \setminus A) \leq \mu(A) + \mu(B)$.
\end{proof}

\section{Jordan Content on $\bb{R}^n$}

\begin{definition}[Rectangle]
A \textbf{rectangle} (or \textbf{box}) in $\bb{R}^n$ is a set of the form:
\[ R = [a_1, b_1] \times [a_2, b_2] \times \cdots \times [a_n, b_n] \]
The \textbf{volume} of $R$ is $\text{vol}(R) = \prod_{i=1}^n (b_i - a_i)$.
\end{definition}

\begin{definition}[Elementary Set]
An \textbf{elementary set} is a finite union of rectangles. The collection of elementary sets forms an algebra, denoted $\mathcal{E}$.
\end{definition}

\begin{proposition}
Every elementary set can be written as a finite disjoint union of rectangles.
\end{proposition}

\begin{proof}
Proceed by induction on the number of rectangles in the union. Base case is trivial.

For the inductive step, let $E = R \cup E'$ where $E'$ is a disjoint union of rectangles $R_1, \ldots, R_k$.

Decompose $R$ into disjoint pieces: $R = (R \setminus \bigcup R_i) \sqcup \bigsqcup (R \cap R_i)$.

Each $R \cap R_i$ is a rectangle (intersection of rectangles). The set $R \setminus \bigcup R_i$ can be written as a finite disjoint union of rectangles by slicing along the faces of each $R_i$.
\end{proof}

\begin{definition}[Jordan Content]
The \textbf{Jordan content} $c$ on $\mathcal{E}$ is defined by:
\[ c(E) = \sum_{i=1}^k \text{vol}(R_i) \]
where $E = \bigsqcup_{i=1}^k R_i$ is a disjoint decomposition into rectangles.
\end{definition}

\begin{proposition}
Jordan content is well-defined and is a content on $\mathcal{E}$.
\end{proposition}

\begin{proof}
\textbf{Well-defined:} If $E = \bigsqcup R_i = \bigsqcup S_j$ are two disjoint decompositions, then $R_i = \bigsqcup_j (R_i \cap S_j)$. Since volume is additive on disjoint rectangles:
\[ \sum \text{vol}(R_i) = \sum_i \sum_j \text{vol}(R_i \cap S_j) = \sum_j \sum_i \text{vol}(R_i \cap S_j) = \sum \text{vol}(S_j) \]

\textbf{Content:} $c(\emptyset) = 0$ (empty sum). For disjoint $E, F \in \mathcal{E}$, write $E = \bigsqcup R_i$, $F = \bigsqcup S_j$. Then $E \cup F = (\bigsqcup R_i) \sqcup (\bigsqcup S_j)$, so $c(E \cup F) = c(E) + c(F)$.
\end{proof}

\section{Inner and Outer Content}

\begin{definition}[Outer Content]
For any bounded $A \subseteq \bb{R}^n$, the \textbf{outer content} is:
\[ \JordanOuter{A} = \inf \{ \Jordan{}(E) \mid E \in \mathcal{E}, A \subseteq E \} \]
\end{definition}

\begin{definition}[Inner Content]
For any bounded $A \subseteq \bb{R}^n$, the \textbf{inner content} is:
\[ \JordanInner{}(A) = \sup \{ \Jordan{}(E) \mid E \in \mathcal{E}, E \subseteq A \} \]
\end{definition}

\begin{proposition}[Properties]
For bounded sets $A, B \subseteq \bb{R}^n$:
\begin{enumerate}
    \item $\JordanInner{}(A) \leq \JordanOuter{}(A)$
    \item $A \subseteq B \implies \JordanInner{}(A) \leq \JordanInner{}(B)$ and $\JordanOuter{}(A) \leq \JordanOuter{}(B)$
    \item $\JordanOuter{}(A \cup B) \leq \JordanOuter{}(A) + \JordanOuter{}(B)$
\end{enumerate}
\end{proposition}

\begin{proof}
(1) For any $E_1 \subseteq A \subseteq E_2$ with $E_1, E_2 \in \mathcal{E}$, we have $\Jordan{E_1} \leq \Jordan{E_2}$. Taking supremum over $E_1$ and infimum over $E_2$ gives $\JordanInner{}(A) \leq \JordanOuter{}(A)$.

(2) Monotonicity of inner/outer content follows from set inclusion.

(3) If $A \subseteq E_1$ and $B \subseteq E_2$, then $A \cup B \subseteq E_1 \cup E_2$ and $\Jordan{E_1 \cup E_2} \leq \Jordan{E_1} + \Jordan{E_2}$.
\end{proof}

\begin{definition}[Jordan Measurable]
A bounded set $A$ is \textbf{Jordan measurable} if $\JordanInner{}(A) = \JordanOuter{}(A)$. The common value is the \textbf{Jordan content} $\Jordan{}(A)$.
\end{definition}

\begin{theorem}[Characterization of Jordan Measurability]
A bounded set $A$ is Jordan measurable if and only if for every $\epsilon > 0$, there exist elementary sets $E_1 \subseteq A \subseteq E_2$ with $\Jordan{E_2 \setminus E_1} < \epsilon$.
\end{theorem}

\begin{proof}
($\Rightarrow$) Given $\epsilon > 0$, by definition of sup and inf, there exist $E_1 \subseteq A$ with $\Jordan{E_1} > \JordanInner{}(A) - \epsilon/2$ and $E_2 \supseteq A$ with $\Jordan{E_2} < \JordanOuter{}(A) + \epsilon/2$. Since $\JordanInner{}(A) = \JordanOuter{}(A)$:
\[ \Jordan{E_2 \setminus E_1} = \Jordan{E_2} - \Jordan{E_1} < \epsilon \]

($\Leftarrow$) We have $\JordanInner{}(A) \geq \Jordan{E_1}$ and $\JordanOuter{}(A) \leq \Jordan{E_2}$. Thus:
\[ \JordanOuter{}(A) - \JordanInner{}(A) \leq \Jordan{E_2} - \Jordan{E_1} = \Jordan{E_2 \setminus E_1} < \epsilon \]
Since $\epsilon$ is arbitrary, $\JordanInner{}(A) = \JordanOuter{}(A)$.
\end{proof}

\begin{theorem}[Boundary Characterization]
A bounded set $A$ is Jordan measurable if and only if $\JordanOuter{\Boundary{A}} = 0$.
\end{theorem}

\begin{proof}
Note that $\Boundary{A} = \Closure{A} \setminus \Interior{A}$, where $\Closure{A}$ is the closure and $\Interior{A}$ is the interior.

($\Rightarrow$) Assume $A$ is Jordan measurable. For any $\epsilon > 0$, there exist elementary sets $E_1 \subseteq A \subseteq E_2$ with $\Jordan{E_2 \setminus E_1} < \epsilon$.

Since $E_1 \subseteq A$, we have $\Interior{E_1} \subseteq \Interior{A}$. Since $A \subseteq E_2$, we have $\Closure{A} \subseteq \Closure{E_2}$.

Thus $\Boundary{A} \subseteq \Closure{E_2} \setminus \Interior{E_1} \subseteq E_2 \setminus E_1$ (for appropriate elementary sets).

So $\JordanOuter{\Boundary{A}} \leq \Jordan{E_2 \setminus E_1} < \epsilon$. Since $\epsilon$ is arbitrary, $\JordanOuter{\Boundary{A}} = 0$.

($\Leftarrow$) Assume $\JordanOuter{\Boundary{A}} = 0$. For any $\epsilon > 0$, cover $\Boundary{A}$ by an elementary set $E$ with $\Jordan{E} < \epsilon$.

Take $E_1 = \Closure{\Interior{A} \setminus E}$ (elementary approximation of interior) and $E_2 = \Closure{A} \cup E$.

Then $E_1 \subseteq A \subseteq E_2$ and $\Jordan{E_2 \setminus E_1} \leq \Jordan{E} < \epsilon$.
\end{proof}


\chapter{The Riemann Integral}

\section{Riemann/Jordan Simple Functions}

\begin{definition}[Riemann Simple Function]
A function $s: R \to \bb{R}$ defined on a rectangle $R \subseteq \bb{R}^n$ is \textbf{Riemann simple} (or \textbf{step function}) if:
\begin{enumerate}
    \item $s$ has finite image
    \item For any subset $B \subseteq \bb{R}$, the preimage $s^{-1}(B)$ is a finite disjoint union of rectangles
\end{enumerate}
Equivalently, $s = \sum_{i=1}^k \alpha_i \CharFunc{R_i}$ where $R_i$ are disjoint rectangles.
\end{definition}

\begin{definition}[Integral of Riemann Simple Function]
For a Riemann simple function $s = \sum_{i=1}^k \alpha_i \CharFunc{R_i}$:
\[ \int_R s \, dc = \sum_{i=1}^k \alpha_i \cdot c(R_i) \]
\end{definition}

\section{Upper and Lower Riemann Sums}

\begin{definition}[Partition]
A \textbf{partition} $P$ of a rectangle $R = [a_1, b_1] \times \cdots \times [a_n, b_n]$ is a collection of subrectangles $\{R_j\}$ such that:
\begin{enumerate}
    \item $R = \bigcup_j R_j$
    \item The interiors of distinct $R_j$ are disjoint
\end{enumerate}
\end{definition}

\begin{definition}[Upper and Lower Sums]
For a bounded function $f: R \to \bb{R}$ and partition $P = \{R_j\}$:
\begin{align*}
    \LowerSum{f}{P} &= \sum_j \left(\inf_{x \in R_j} f(x)\right) \cdot c(R_j) \\
    \UpperSum{f}{P} &= \sum_j \left(\sup_{x \in R_j} f(x)\right) \cdot c(R_j)
\end{align*}
\end{definition}

\begin{proposition}
For any partitions $P, Q$ of $R$:
\[ \LowerSum{f}{P} \leq \UpperSum{f}{Q} \]
\end{proposition}

\begin{proof}
Let $P \vee Q$ be the common refinement. Then:
\[ \LowerSum{f}{P} \leq \LowerSum{f}{P \vee Q} \leq \UpperSum{f}{P \vee Q} \leq \UpperSum{f}{Q} \]
\end{proof}

\section{Riemann Integral}

\begin{definition}[Upper and Lower Riemann Integrals]
\begin{align*}
    \LowerRiemann_R f \, dc &= \sup_P \LowerSum{f}{P} \\
    \UpperRiemann_R f \, dc &= \inf_P \UpperSum{f}{P}
\end{align*}
\end{definition}

\begin{definition}[Riemann Integrable]
$f$ is \textbf{Riemann integrable} on $R$ if $\LowerRiemann_R f = \UpperRiemann_R f$. The common value is $\int_R f \, dc$.
\end{definition}

\begin{theorem}[Characterization of Riemann Integrability]
$f$ is Riemann integrable if and only if for every $\epsilon > 0$, there exists a partition $P$ with:
\[ \UpperSum{f}{P} - \LowerSum{f}{P} < \epsilon \]
\end{theorem}

\begin{proof}
($\Rightarrow$) Given $\epsilon > 0$, choose partitions $P_1, P_2$ with $\LowerSum{f}{P_1} > \LowerRiemann f - \epsilon/2$ and $\UpperSum{f}{P_2} < \UpperRiemann f + \epsilon/2$. For $P = P_1 \vee P_2$:
\[ \UpperSum{f}{P} - \LowerSum{f}{P} \leq \UpperSum{f}{P_2} - \LowerSum{f}{P_1} < \UpperRiemann f - \LowerRiemann f + \epsilon = \epsilon \]

($\Leftarrow$) For any $\epsilon > 0$:
\[ \UpperRiemann f - \LowerRiemann f \leq \UpperSum{f}{P} - \LowerSum{f}{P} < \epsilon \]
\end{proof}

\section{Riemann Double Integral}

\begin{definition}[Iterated Integral]
For $f: R_1 \times R_2 \to \bb{R}$ where $R_i \subseteq \bb{R}^{n_i}$, the \textbf{iterated integral} is:
\[ \int_{R_1} \left( \int_{R_2} f(x, y) \, dc(y) \right) dc(x) \]
when both inner integrals exist.
\end{definition}

\begin{theorem}[Fubini for Riemann Integral]
If $f: R_1 \times R_2 \to \bb{R}$ is Riemann integrable, and for each $x \in R_1$ the function $y \mapsto f(x, y)$ is Riemann integrable on $R_2$, then:
\[ \int_{R_1 \times R_2} f \, dc = \int_{R_1} \left( \int_{R_2} f(x, y) \, dc(y) \right) dc(x) \]
\end{theorem}

\begin{proof}
Let $P = P_1 \times P_2$ be a partition of $R_1 \times R_2$ with $P_i$ partitions of $R_i$.

For each subrectangle $S_1 \times S_2$ in $P$:
\[ \inf_{S_1 \times S_2} f \leq \inf_{x \in S_1} \left( \inf_{y \in S_2} f(x, y) \right) \leq \inf_{x \in S_1} \int_{S_2} f(x, y) dc(y) / c(S_2) \]

Summing over subrectangles:
\[ \LowerSum{f}{P} \leq \LowerSum{g}{P_1} \]
where $g(x) = \int_{R_2} f(x, y) dc(y)$.

Similarly $\UpperSum{g}{P_1} \leq \UpperSum{f}{P}$.

Since $f$ is Riemann integrable, $\UpperSum{f}{P} - \LowerSum{f}{P} \to 0$, forcing $g$ to be Riemann integrable with the stated equality.
\end{proof}

\section{Limitations of Riemann Integration}

\begin{example}[Dirichlet Function]
The function $\CharFunc{\bb{Q} \cap [0,1]}$ is not Riemann integrable. For any partition:
\[ \underline{S} = 0, \quad \overline{S} = 1 \]
since every interval contains both rationals and irrationals.
\end{example}

\begin{remark}
Riemann integration has fundamental limitations:
\begin{enumerate}
    \item Cannot integrate many discontinuous functions
    \item Poor behavior under limits (no good convergence theorems)
    \item Requires bounded functions on bounded domains
\end{enumerate}
These motivate the development of Lebesgue integration.
\end{remark}


\chapter{Measures}

\section{Definition of Measure}

\begin{definition}[Measure]
Let $(X, \mathcal{M})$ be a measurable space. A function $\mu: \mathcal{M} \to [0, \infty]$ is a \textbf{measure} if:
\begin{enumerate}
    \item $\mu(\emptyset) = 0$
    \item \textbf{$\sigma$-additivity:} For any sequence of pairwise disjoint sets $(A_n)_{n \in \bb{N}} \subseteq \mathcal{M}$:
    \[ \mu\left(\bigcup_{n=1}^{\infty} A_n\right) = \sum_{n=1}^{\infty} \mu(A_n) \]
\end{enumerate}
The triple $(X, \mathcal{M}, \mu)$ is called a \textbf{measure space}.
\end{definition}

\begin{proposition}[Properties of Measures]
Let $\mu$ be a measure on $(X, \mathcal{M})$. Then:
\begin{enumerate}
    \item \textbf{Finite additivity:} $\mu$ is a content
    \item \textbf{Monotonicity:} $A \subseteq B \implies \mu(A) \leq \mu(B)$
    \item \textbf{$\sigma$-subadditivity:} $\mu(\bigcup_{n=1}^\infty A_n) \leq \sum_{n=1}^\infty \mu(A_n)$
    \item \textbf{Continuity from below:} If $A_n \nearrow A$, then $\mu(A) = \lim_{n \to \infty} \mu(A_n)$
    \item \textbf{Continuity from above:} If $A_n \searrow A$ and $\mu(A_k) < \infty$ for some $k$, then $\mu(A) = \lim_{n \to \infty} \mu(A_n)$
\end{enumerate}
\end{proposition}

\begin{proof}
(1) Take $A_3 = A_4 = \cdots = \emptyset$.

(2) $\mu(B) = \mu(A) + \mu(B \setminus A) \geq \mu(A)$.

(3) Define $B_1 = A_1$, $B_n = A_n \setminus \bigcup_{i=1}^{n-1} A_i$. Then $B_n$ are disjoint, $\bigcup A_n = \bigcup B_n$, and $B_n \subseteq A_n$. Thus:
\[ \mu\left(\bigcup A_n\right) = \mu\left(\bigcup B_n\right) = \sum \mu(B_n) \leq \sum \mu(A_n) \]

(4) Define $B_1 = A_1$, $B_n = A_n \setminus A_{n-1}$. Then $B_n$ are disjoint and $A_n = \bigcup_{i=1}^n B_i$. Since $A = \bigcup A_n = \bigcup B_n$:
\[ \mu(A) = \sum_{n=1}^\infty \mu(B_n) = \lim_{N \to \infty} \sum_{n=1}^N \mu(B_n) = \lim_{N \to \infty} \mu(A_N) \]

(5) WLOG assume $\mu(A_k) < \infty$ for some $k$ (by hypothesis). Define $C_n = A_k \setminus A_n$ for $n \geq k$. Then $C_n \nearrow A_k \setminus A$, so:
\[ \mu(A_k \setminus A) = \lim_{n \to \infty} \mu(C_n) = \lim_{n \to \infty} (\mu(A_k) - \mu(A_n)) \]
Thus $\mu(A) = \mu(A_k) - \mu(A_k \setminus A) = \lim \mu(A_n)$.
\end{proof}

\begin{example}
The finiteness condition in continuity from above is essential. Let $A_n = [n, \infty) \subseteq \bb{R}$ with Lebesgue measure. Then $A_n \searrow \emptyset$ but $\lambda(A_n) = \infty$ for all $n$.
\end{example}

\section{Outer Measures}

\begin{definition}[Outer Measure]
A function $\mu^*: \Powerset{X} \to [0, \infty]$ is an \textbf{outer measure} if:
\begin{enumerate}
    \item $\mu^*(\emptyset) = 0$
    \item \textbf{Monotonicity:} $A \subseteq B \implies \mu^*(A) \leq \mu^*(B)$
    \item \textbf{$\sigma$-subadditivity:} $\mu^*(\bigcup_{n=1}^\infty A_n) \leq \sum_{n=1}^\infty \mu^*(A_n)$
\end{enumerate}
\end{definition}

\begin{definition}[Lebesgue Outer Measure]
For $A \subseteq \bb{R}^n$:
\[ \lambda^*(A) = \inf \left\{ \sum_{k=1}^\infty \text{vol}(R_k) \mid A \subseteq \bigcup_{k=1}^\infty R_k, R_k \text{ open rectangles} \right\} \]
\end{definition}

\begin{proposition}
$\lambda^*$ is an outer measure on $\bb{R}^n$.
\end{proposition}

\begin{proof}
\textbf{(1) $\lambda^*(\emptyset) = 0$:} Cover $\emptyset$ by rectangles of arbitrarily small total volume.

\textbf{(2) Monotonicity:} If $A \subseteq B$, any cover of $B$ is a cover of $A$, so $\lambda^*(A) \leq \lambda^*(B)$.

\textbf{(3) $\sigma$-subadditivity:} Let $(A_n)$ be a sequence. For each $n$, given $\epsilon > 0$, choose open rectangles $\{R_{n,k}\}_k$ covering $A_n$ with $\sum_k \text{vol}(R_{n,k}) < \lambda^*(A_n) + \epsilon/2^n$.

Then $\{R_{n,k}\}_{n,k}$ covers $\bigcup A_n$ and:
\[ \sum_{n,k} \text{vol}(R_{n,k}) < \sum_n (\lambda^*(A_n) + \epsilon/2^n) = \sum_n \lambda^*(A_n) + \epsilon \]

Thus $\lambda^*(\bigcup A_n) \leq \sum \lambda^*(A_n) + \epsilon$. Since $\epsilon$ is arbitrary, done.
\end{proof}

\section{Complete Measures and Completion}

\begin{definition}[Complete Measure]
A measure space $(X, \mathcal{M}, \mu)$ is \textbf{complete} if every subset of a null set is measurable:
\[ \mu(N) = 0 \text{ and } E \subseteq N \implies E \in \mathcal{M} \]
\end{definition}

\begin{definition}[Completion of a Measure Space]
Let $(X, \mathcal{M}, \mu)$ be a measure space. Its \textbf{completion} is the measure space $(X, \overline{\mathcal{M}}, \overline{\mu})$ where:
\[ \overline{\mathcal{M}} = \{ A \cup N \mid A \in \mathcal{M}, \Exists{M \in \mathcal{M}} N \subseteq M \text{ and } \mu(M) = 0 \} \]
and $\overline{\mu}(A \cup N) = \mu(A)$.
\end{definition}

\begin{proposition}
The completion $(X, \overline{\mathcal{M}}, \overline{\mu})$ is a complete measure space.
\end{proposition}

\begin{proof}
Let $E \subseteq N'$ where $\overline{\mu}(N') = 0$. Then $N' = A \cup N$ with $\mu(A) = 0$ and $N \subseteq M$ for some $M \in \mathcal{M}$ with $\mu(M) = 0$.

Since $E \subseteq N' = A \cup N \subseteq A \cup M$, and $\mu(A \cup M) = 0$, we have $E \in \overline{\mathcal{M}}$ (taking $A' = \emptyset$ and $N'' = E$).
\end{proof}

\section{Carathéodory's Theorem}

\begin{definition}[Carathéodory Measurability]
A set $E \subseteq X$ is \textbf{$\mu^*$-measurable} if for every $A \subseteq X$:
\[ \mu^*(A) = \mu^*(A \cap E) + \mu^*(A \cap E^c) \]
\end{definition}

\begin{theorem}[Carathéodory]
Let $\mu^*$ be an outer measure on $X$. Let $\mathcal{M}(\mu^*)$ be the collection of $\mu^*$-measurable sets. Then:
\begin{enumerate}
    \item $\mathcal{M}(\mu^*)$ is a $\sigma$-algebra
    \item $\mu^*|_{\mathcal{M}(\mu^*)}$ is a complete measure
\end{enumerate}
\end{theorem}

\begin{proof}
\textbf{Step 1 (Complements):} $\mathcal{M}(\mu^*)$ is closed under complements by symmetry of the definition: $E$ is measurable iff $E^c$ is measurable.

\textbf{Step 2 (Finite unions):} Let $E, F \in \mathcal{M}(\mu^*)$. We show $E \cup F \in \mathcal{M}(\mu^*)$.

For any test set $A \subseteq X$, we need $\mu^*(A) = \mu^*(A \cap (E \cup F)) + \mu^*(A \cap (E \cup F)^c)$.

Since $E$ is measurable: $\mu^*(A) = \mu^*(A \cap E) + \mu^*(A \cap E^c)$.

Since $F$ is measurable applied to $A \cap E^c$:
\[ \mu^*(A \cap E^c) = \mu^*(A \cap E^c \cap F) + \mu^*(A \cap E^c \cap F^c) \]

Note that $(E \cup F)^c = E^c \cap F^c$, so:
\[ \mu^*(A) = \mu^*(A \cap E) + \mu^*(A \cap E^c \cap F) + \mu^*(A \cap (E \cup F)^c) \]

By subadditivity: $\mu^*(A \cap E) + \mu^*(A \cap E^c \cap F) \geq \mu^*(A \cap (E \cup F))$.

Thus $\mu^*(A) \geq \mu^*(A \cap (E \cup F)) + \mu^*(A \cap (E \cup F)^c)$.

The reverse inequality holds by subadditivity. Hence $E \cup F \in \mathcal{M}(\mu^*)$.

\textbf{Step 3 (Countable unions):} Let $(E_n)$ be disjoint sets in $\mathcal{M}(\mu^*)$. Let $F_n = \bigcup_{i=1}^n E_i$.

\textit{Claim:} For any $A$, $\mu^*(A \cap F_n) = \sum_{i=1}^n \mu^*(A \cap E_i)$.

\textit{Proof by induction:} Base case $n=1$ is trivial. Assume true for $n$. Since $E_{n+1}$ is measurable:
\[ \mu^*(A \cap F_{n+1}) = \mu^*(A \cap F_{n+1} \cap E_{n+1}) + \mu^*(A \cap F_{n+1} \cap E_{n+1}^c) \]
Since $E_{n+1} \subseteq F_{n+1}$ and $F_n = F_{n+1} \cap E_{n+1}^c$:
\[ \mu^*(A \cap F_{n+1}) = \mu^*(A \cap E_{n+1}) + \mu^*(A \cap F_n) = \mu^*(A \cap E_{n+1}) + \sum_{i=1}^n \mu^*(A \cap E_i) \]

Now let $E = \bigcup_{n=1}^\infty E_n$. Since $F_n \in \mathcal{M}(\mu^*)$ and $F_n \subseteq E$:
\[ \mu^*(A) = \mu^*(A \cap F_n) + \mu^*(A \cap F_n^c) \geq \sum_{i=1}^n \mu^*(A \cap E_i) + \mu^*(A \cap E^c) \]

Taking $n \to \infty$: $\mu^*(A) \geq \sum_{i=1}^\infty \mu^*(A \cap E_i) + \mu^*(A \cap E^c) \geq \mu^*(A \cap E) + \mu^*(A \cap E^c)$.

The reverse inequality holds by subadditivity. Thus $E \in \mathcal{M}(\mu^*)$ and $\mu^*$ is $\sigma$-additive on $\mathcal{M}(\mu^*)$.

\textbf{Step 4 (Completeness):} Let $\mu^*(N) = 0$ and $E \subseteq N$. For any $A$:
\[ \mu^*(A \cap E) \leq \mu^*(N) = 0 \]
\[ \mu^*(A) \leq \mu^*(A \cap E) + \mu^*(A \cap E^c) = \mu^*(A \cap E^c) \leq \mu^*(A) \]
Thus $\mu^*(A) = \mu^*(A \cap E) + \mu^*(A \cap E^c)$, so $E \in \mathcal{M}(\mu^*)$.
\end{proof}

\section{Lebesgue Measure}

\begin{definition}
The \textbf{Lebesgue measure} $\lambda$ on $\bb{R}^n$ is the restriction of $\lambda^*$ to $\mathcal{M}(\lambda^*)$, the Lebesgue measurable sets.
\end{definition}

\begin{theorem}
Every Borel set is Lebesgue measurable: $\BorelAlg{\bb{R}^n} \subseteq \mathcal{M}(\lambda^*)$.

For any rectangle $R$: $\lambda(R) = \text{vol}(R)$.
\end{theorem}

\begin{proof}
\textbf{Step 1:} Open sets are measurable. Let $G$ be open. For any test set $A$:

Since $\lambda^*$ is defined using open covers, for any $\epsilon > 0$ there exist open rectangles $\{R_k\}$ with $A \subseteq \bigcup R_k$ and $\sum \text{vol}(R_k) < \lambda^*(A) + \epsilon$.

Using the fact that $G$ is open to split each $R_k$ appropriately, one shows:
\[ \lambda^*(A \cap G) + \lambda^*(A \cap G^c) \leq \lambda^*(A) + \epsilon \]

Since $\epsilon$ is arbitrary, $G \in \mathcal{M}(\lambda^*)$.

\textbf{Step 2:} Since $\mathcal{M}(\lambda^*)$ is a $\sigma$-algebra containing all open sets, and $\BorelAlg{\bb{R}^n} = \GenSigmaAlg{\text{open sets}}$, we have $\BorelAlg{\bb{R}^n} \subseteq \mathcal{M}(\lambda^*)$.

\textbf{Step 3:} For rectangles, $\lambda^*(R) \leq \text{vol}(R)$ by covering with $R$ itself. The reverse follows from $\sigma$-subadditivity and convergence arguments.
\end{proof}

\section{Regularity}

\begin{theorem}[Outer Regularity]
For any $A \subseteq \bb{R}^n$:
\[ \lambda^*(A) = \inf \{ \lambda(G) \mid G \text{ open}, G \supseteq A \} \]
\end{theorem}

\begin{proof}
Let $\alpha = \inf \{ \lambda(G) \mid G \text{ open}, G \supseteq A \}$.

Since $\lambda^*(A) \leq \lambda(G)$ for any open $G \supseteq A$, we have $\lambda^*(A) \leq \alpha$.

For the reverse: given $\epsilon > 0$, by definition of $\lambda^*$, there exist open rectangles $\{R_k\}$ with $A \subseteq \bigcup R_k$ and $\sum \text{vol}(R_k) < \lambda^*(A) + \epsilon$.

Let $G = \bigcup R_k$, which is open. Then $\lambda(G) \leq \sum \text{vol}(R_k) < \lambda^*(A) + \epsilon$.

Thus $\alpha \leq \lambda^*(A) + \epsilon$, so $\alpha = \lambda^*(A)$.
\end{proof}

\begin{theorem}[Inner Regularity]
For any Lebesgue measurable $A$:
\[ \lambda(A) = \sup \{ \lambda(K) \mid K \text{ compact}, K \subseteq A \} \]
\end{theorem}

\begin{proof}
\textbf{Case 1:} $A$ bounded. Let $A \subseteq [-M, M]^n$.

By outer regularity of $A^c \cap [-M,M]^n$, for any $\epsilon > 0$ there exists open $G \supseteq A^c \cap [-M,M]^n$ with $\lambda(G) < \lambda(A^c \cap [-M,M]^n) + \epsilon$.

Let $K = [-M,M]^n \setminus G$. Then $K$ is closed and bounded (hence compact), $K \subseteq A$, and:
\[ \lambda(K) = \lambda([-M,M]^n) - \lambda(G) > \lambda(A) - \epsilon \]

\textbf{Case 2:} $A$ unbounded. Write $A = \bigcup_n (A \cap [-n,n]^n)$ and use continuity from below.
\end{proof}

\section{Completion}

\begin{theorem}
The Lebesgue measure is the completion of the Borel measure on $\bb{R}^n$.

Equivalently: $A \in \mathcal{M}(\lambda^*)$ iff $A = B \cup N$ where $B \in \BorelAlg{\bb{R}^n}$ and $\lambda^*(N) = 0$.
\end{theorem}

\begin{proof}
\textbf{($\Leftarrow$)} If $B \in \BorelAlg{\bb{R}^n}$ and $\lambda^*(N) = 0$, then $B$ is measurable (Borel $\subseteq$ Lebesgue) and $N$ is measurable by completeness of $\lambda$. Thus $B \cup N$ is measurable.

\textbf{($\Rightarrow$)} Let $A \in \mathcal{M}(\lambda^*)$. By outer regularity, for each $n$ there exists open $G_n \supseteq A$ with $\lambda(G_n) < \lambda(A) + 1/n$.

Let $B = \bigcap_n G_n$. Then $B \in \BorelAlg{\bb{R}^n}$, $A \subseteq B$, and $\lambda(B) = \lambda(A)$.

Let $N = B \setminus A$. Then $\lambda(N) = \lambda(B) - \lambda(A) = 0$ and $A = B \setminus N$.

Since $A \subseteq B$, we can write $A = (B \setminus N) = B \cap N^c$. Rearranging: $B = A \cup N$ where $\lambda^*(N) = 0$.
\end{proof}

\section{Uniqueness}

\begin{theorem}[Uniqueness via $\pi$-Systems]
Let $\mu, \nu$ be $\sigma$-finite measures on $(X, \mathcal{M})$. If $\mathcal{P}$ is a $\pi$-system with $\GenSigmaAlg{\mathcal{P}} = \mathcal{M}$ and $\mu|_{\mathcal{P}} = \nu|_{\mathcal{P}}$, then $\mu = \nu$.
\end{theorem}

\begin{proof}
Define $\mathcal{L} = \{A \in \mathcal{M} \mid \mu(A) = \nu(A)\}$. We show $\mathcal{L}$ is a $\lambda$-system.

\textbf{(i)} $X \in \mathcal{L}$: By $\sigma$-finiteness, write $X = \bigcup_{n=1}^\infty X_n$ with $X_n \nearrow X$ and $\mu(X_n), \nu(X_n) < \infty$. If $\mu|_{\mathcal{P}} = \nu|_{\mathcal{P}}$ and $X_n \in \mathcal{P}$, then $\mu(X_n) = \nu(X_n)$, so $\mu(X) = \nu(X)$ by continuity.

\textbf{(ii)} If $A \subseteq B$ with $A, B \in \mathcal{L}$ and $\mu(B) < \infty$:
\[ \mu(B \setminus A) = \mu(B) - \mu(A) = \nu(B) - \nu(A) = \nu(B \setminus A) \]
So $B \setminus A \in \mathcal{L}$.

For the general case with $\sigma$-finiteness: intersect with $X_n$ and use continuity.

\textbf{(iii)} If $A_n \nearrow A$ with $A_n \in \mathcal{L}$:
\[ \mu(A) = \lim_n \mu(A_n) = \lim_n \nu(A_n) = \nu(A) \]

Since $\mathcal{P} \subseteq \mathcal{L}$ and $\mathcal{L}$ is a $\lambda$-system, by Dynkin's theorem: $\mathcal{M} = \GenSigmaAlg{\mathcal{P}} = \GenLambda{\mathcal{P}} \subseteq \mathcal{L}$.
\end{proof}


\chapter{The Lebesgue Integral}

\section{Measurable Functions}

\begin{definition}[Measurable Function]
Let $(X, \mathcal{M})$ and $(Y, \mathcal{N})$ be measurable spaces. A function $f: X \to Y$ is \textbf{measurable} if:
\[ \Forall :{A \in \mathcal{N}}{f^{-1}(A) \in \mathcal{M}} \]
\end{definition}

\begin{proposition}[Characterization for Real Functions]
$f: X \to \bb{R}$ is measurable iff any of the following equivalent conditions hold:
\begin{enumerate}
    \item $f^{-1}((a, \infty)) \in \mathcal{M}$ for all $a \in \bb{R}$
    \item $f^{-1}([a, \infty)) \in \mathcal{M}$ for all $a \in \bb{R}$
    \item $f^{-1}((-\infty, a)) \in \mathcal{M}$ for all $a \in \bb{R}$
    \item $f^{-1}((-\infty, a]) \in \mathcal{M}$ for all $a \in \bb{R}$
\end{enumerate}
\end{proposition}

\begin{proof}
These four families generate $\BorelAlg{\bb{R}}$. Thus any condition implies $f^{-1}(B) \in \mathcal{M}$ for all $B \in \BorelAlg{\bb{R}}$.

The equivalences follow from: $(a, \infty)^c = (-\infty, a]$, $[a, \infty) = \bigcap_n (a - 1/n, \infty)$, etc.
\end{proof}

\begin{proposition}[Closure Properties]
If $f, g$ are measurable and $c \in \bb{R}$:
\begin{enumerate}
    \item $cf$, $f + g$, $fg$, $|f|$, $\max(f, g)$, $\min(f, g)$ are measurable
    \item If $(f_n)$ are measurable, then $\sup_n f_n$, $\inf_n f_n$, $\limsup f_n$, $\liminf f_n$ are measurable
\end{enumerate}
\end{proposition}

\begin{proof}
(1) For $cf$: $(cf)^{-1}((a, \infty)) = f^{-1}((a/c, \infty))$ (case $c > 0$).

For $f + g$: $(f+g)^{-1}((a, \infty)) = \bigcup_{r \in \bb{Q}} (f^{-1}((r, \infty)) \cap g^{-1}((a-r, \infty)))$.

For $fg$: use $fg = \frac{1}{4}((f+g)^2 - (f-g)^2)$ and $h^2$ is measurable since $(h^2)^{-1}((a, \infty)) = h^{-1}((\sqrt{a}, \infty)) \cup h^{-1}((-\infty, -\sqrt{a}))$ for $a \geq 0$.

For $\max$: $(\max(f,g))^{-1}((a, \infty)) = f^{-1}((a, \infty)) \cup g^{-1}((a, \infty))$.

(2) $(\sup_n f_n)^{-1}((a, \infty)) = \bigcup_n f_n^{-1}((a, \infty)) \in \mathcal{M}$. Similarly for $\inf$. Then $\limsup f_n = \inf_k \sup_{n \geq k} f_n$.
\end{proof}

\section{Simple Functions}

\begin{definition}[Simple Function]
A \textbf{simple function} is $s: X \to [0, \infty)$ with finite image. Canonical form:
\[ s = \sum_{i=1}^n \alpha_i \CharFunc{A_i} \]
where $\alpha_i$ are distinct values and $A_i = s^{-1}(\{\alpha_i\})$ are disjoint.
\end{definition}

\begin{definition}[$\sigma$-Simple Function]
A function $s: X \to [0, \infty]$ is \textbf{$\sigma$-simple} if its image is at most countable and each level set is measurable. Canonical form:
\[ s = \sum_{i=1}^\infty \alpha_i \CharFunc{A_i} \]
where $(\alpha_i)_{i \in \bb{N}}$ are distinct values in $[0, \infty]$ and $A_i = s^{-1}(\{\alpha_i\})$ are disjoint measurable sets.
\end{definition}

\section{Simple Function Approximation}

\begin{theorem}[Increasing Simple Function Approximation]
Let $f: X \to [0, \infty]$ be measurable. Then there exists a sequence $(s_n)$ of simple functions such that:
\begin{enumerate}
    \item $0 \leq s_1 \leq s_2 \leq \cdots \leq f$ (increasing)
    \item $s_n \to f$ pointwise
\end{enumerate}
Moreover, $f$ is bounded if and only if the convergence is uniform.
\end{theorem}

\begin{proof}
Define:
\[ s_n(x) = \frac{\lfloor n f(x) \rfloor}{n} \cdot \CharFunc{\{f < n\}}(x) + n \cdot \CharFunc{\{f \geq n\}}(x) \]

Then $s_n$ is simple with finitely many values in $\{0, 1/n, 2/n, \ldots, n\}$.

\textbf{Verification:}
\begin{itemize}
    \item $s_n \leq f$: If $f(x) < n$, then $s_n(x) = \lfloor nf(x) \rfloor / n \leq f(x)$. If $f(x) \geq n$, then $s_n(x) = n \leq f(x)$.
    \item $s_n$ increasing: As $n$ increases, the floor function gives finer approximation and the cap $n$ increases.
    \item $s_n \to f$: If $f(x) < \infty$, for $n > f(x)$: $f(x) - s_n(x) < 1/n \to 0$. If $f(x) = \infty$, $s_n(x) = n \to \infty$.
\end{itemize}

\textbf{Uniform iff bounded:}
\begin{itemize}
    \item ($\Rightarrow$) Uniform convergence implies $f < s_n + 1 \leq n + 1$ for large $n$, so $f$ is bounded.
    \item ($\Leftarrow$) If $f \leq M$, then for $n > M$: $\|f - s_n\|_\infty < 1/n \to 0$.
\end{itemize}
\end{proof}

\begin{proposition}[Decreasing Simple Function Approximation]
Let $f: X \to [0, \infty]$ be measurable. Then $f$ can be approximated by a decreasing sequence of simple functions $t_n \searrow f$ if and only if $f$ is bounded. In this case, the convergence is uniform.
\end{proposition}

\begin{proof}
($\Leftarrow$) Suppose $f \leq M$. Define:
\[ t_n(x) = \frac{\lceil n f(x) \rceil}{n} \]

Then $t_n$ is simple with finitely many values in $\{0, 1/n, 2/n, \ldots, \lceil nM \rceil / n\}$.

We have $t_n \geq f$, $t_n$ decreases, and $\|t_n - f\|_\infty \leq 1/n \to 0$.

($\Rightarrow$) If $t_n \searrow f$ with $t_n$ simple, then each $t_n$ is bounded. Since $t_1 \geq f$, we have $f \leq \|t_1\|_\infty < \infty$.
\end{proof}

\section{$\sigma$-Simple Function Approximation}

\begin{theorem}[$\sigma$-Simple Function Approximation]
Let $f: X \to [0, \infty]$ be measurable. Then:
\begin{enumerate}
    \item There exists an increasing sequence of $\sigma$-simple functions $s_n \nearrow f$
    \item There exists a decreasing sequence of $\sigma$-simple functions $t_n \searrow f$
\end{enumerate}
Both convergences are always uniform (in the extended metric on $[0, \infty]$).
\end{theorem}

\begin{proof}
\textbf{Increasing:} Define:
\[ s_n(x) = \begin{cases}
\lfloor n f(x) \rfloor / n & \text{if } f(x) < \infty \\
\infty & \text{if } f(x) = \infty
\end{cases} \]

Then $s_n$ is $\sigma$-simple with values in $\{k/n \mid k \in \bb{N}_0\} \cup \{\infty\}$.

Verification: $s_n \leq f$ always. For $f(x) < \infty$: $f(x) - s_n(x) < 1/n$. For $f(x) = \infty$: $s_n(x) = \infty = f(x)$.

The sequence increases since $\lfloor nf \rfloor / n \leq \lfloor (n+1)f \rfloor / (n+1)$ follows from monotonicity of the floor approximation.

\textbf{Decreasing:} Define:
\[ t_n(x) = \begin{cases}
\lceil n f(x) \rceil / n & \text{if } f(x) < \infty \\
\infty & \text{if } f(x) = \infty
\end{cases} \]

Then $t_n$ is $\sigma$-simple with values in $\{k/n \mid k \in \bb{N}_0\} \cup \{\infty\}$.

Verification: $t_n \geq f$ always. For $f(x) < \infty$: $t_n(x) - f(x) < 1/n$. For $f(x) = \infty$: $t_n(x) = \infty = f(x)$.

\textbf{Uniform convergence:} For both sequences, for any $x$:
\[ |s_n(x) - f(x)| < 1/n \quad \text{and} \quad |t_n(x) - f(x)| < 1/n \]
(with the convention that $|\infty - \infty| = 0$). Thus $\|s_n - f\|_\infty, \|t_n - f\|_\infty \leq 1/n \to 0$.
\end{proof}

\section{Integration of Positive Functions}

\begin{definition}[Integral of Simple Function]
For simple $s = \sum_{i=1}^n \alpha_i \CharFunc{A_i}$ on $(X, \mathcal{M}, \mu)$:
\[ \int_X s \, d\mu = \sum_{i=1}^n \alpha_i \mu(A_i) \]
with convention $0 \cdot \infty = 0$.
\end{definition}

\begin{definition}[Integral of $\sigma$-Simple Function]
For $\sigma$-simple $s = \sum_{i=1}^\infty \alpha_i \CharFunc{A_i}$ on $(X, \mathcal{M}, \mu)$:
\[ \int_X s \, d\mu = \sum_{i=1}^\infty \alpha_i \mu(A_i) \]
with convention $0 \cdot \infty = 0$. The sum is always well-defined in $[0, \infty]$.
\end{definition}

\begin{definition}[Lower and Upper Integrals]
For measurable $f: X \to [0, \infty]$, we define four integrals:

\textbf{Lower Simple Integral:}
\[ \underline{\int}_X f \, d\mu = \sup \left\{ \int_X s \, d\mu \mid s \text{ simple}, 0 \leq s \leq f \right\} \]

\textbf{Lower $\sigma$-Simple Integral:}
\[ \underline{\int}^\sigma_X f \, d\mu = \sup \left\{ \int_X s \, d\mu \mid s \text{ $\sigma$-simple}, 0 \leq s \leq f \right\} \]

\textbf{Upper $\sigma$-Simple Integral:}
\[ \overline{\int}^\sigma_X f \, d\mu = \inf \left\{ \int_X t \, d\mu \mid t \text{ $\sigma$-simple}, t \geq f \right\} \]

\textbf{Upper Simple Integral:}
\[ \overline{\int}_X f \, d\mu = \inf \left\{ \int_X t \, d\mu \mid t \text{ simple}, t \geq f \right\} \]
\end{definition}

\begin{proposition}[Ordering]
For any measurable $f: X \to [0, \infty]$:
\[ \underline{\int} f \leq \underline{\int}^\sigma f \leq \overline{\int}^\sigma f \leq \overline{\int} f \]
\end{proposition}

\begin{proof}
The outer inequalities follow since simple $\subseteq$ $\sigma$-simple. The middle inequality follows since $s \leq f \leq t$ implies $\int s \leq \int t$.
\end{proof}

\begin{theorem}[Equality of Integrals]
For any measurable $f: X \to [0, \infty]$:
\[ \underline{\int} f = \underline{\int}^\sigma f = \overline{\int}^\sigma f \]
We denote this common value by $\int f \, d\mu$, the \textbf{Lebesgue integral}.

The upper simple integral $\overline{\int} f$ equals $\int f$ if and only if $f$ is bounded \almosteverywhere{}
\end{theorem}

\begin{proof}
\textbf{Step 1:} $\underline{\int} f = \underline{\int}^\sigma f$.

($\geq$) Immediate since simple $\subseteq$ $\sigma$-simple.

($\leq$) Let $s$ be $\sigma$-simple with $s \leq f$, and let $\epsilon > 0$. By uniform approximation, there exists simple $s'$ with $\|s - s'\|_\infty < \epsilon$. In particular, $s - \epsilon \leq s' \leq s + \epsilon$.

Define $s'' = \max(s' - \epsilon, 0)$, which is simple and satisfies $s'' \leq s \leq f$.

Then $\int s'' \geq \int s - 2\epsilon \cdot \mu(X)$ (if $\mu(X) < \infty$). Taking $\epsilon \to 0$ and supremum over simple functions below $f$ gives $\underline{\int} f \geq \int s$. Taking supremum over $\sigma$-simple $s$: $\underline{\int} f \geq \underline{\int}^\sigma f$.

For $\mu(X) = \infty$, the argument uses truncation to sets of finite measure.

\textbf{Step 2:} $\underline{\int}^\sigma f = \overline{\int}^\sigma f$.

Let $\epsilon > 0$. By uniform $\sigma$-simple approximation, there exist $\sigma$-simple $s, t$ with $s \leq f \leq t$ and $\|t - s\|_\infty < \epsilon$.

Then $\int t - \int s = \int (t - s) \leq \epsilon \cdot \mu(X)$.

Thus $\overline{\int}^\sigma f \leq \int t \leq \int s + \epsilon \cdot \mu(X) \leq \underline{\int}^\sigma f + \epsilon \cdot \mu(X)$.

Since $\epsilon$ is arbitrary (and for $\mu(X) = \infty$, use truncation), we get $\overline{\int}^\sigma f \leq \underline{\int}^\sigma f$.

Combined with $\underline{\int}^\sigma f \leq \overline{\int}^\sigma f$, we have equality.

\textbf{Step 3:} Upper simple integral.

If $f \leq M$ \almosteverywhere{}, then by decreasing simple approximation, there exist simple $t_n \searrow f$. For any such $t_n$: $\overline{\int} f \leq \int t_n$. Since $t_n - f < 1/n$, we get $\int t_n \leq \int f + \mu(X)/n$. Thus $\overline{\int} f = \int f$.

If $f$ is unbounded on a set $A$ of positive measure, then for any simple $t \geq f$, we have $t \geq f$ on $A$. But $t$ is bounded, contradiction. Hence no simple $t \geq f$ exists with $\int t < \infty$, so $\overline{\int} f = \infty$.
\end{proof}

\begin{definition}[Lebesgue Integral]
For measurable $f: X \to [0, \infty]$:
\[ \int_X f \, d\mu = \underline{\int} f = \underline{\int}^\sigma f = \overline{\int}^\sigma f \]
\end{definition}

\section{Monotone Convergence Theorem}

\begin{theorem}[MCT / Beppo-Levi]
Let $(f_n)$ be measurable with $0 \leq f_n \nearrow f$ pointwise. Then:
\[ \lim_{n \to \infty} \int_X f_n \, d\mu = \int_X f \, d\mu \]
\end{theorem}

\begin{proof}
Since $f_n \leq f$, monotonicity gives $\int f_n \leq \int f$, so $\lim \int f_n \leq \int f$.

For the reverse, let $s$ be simple with $0 \leq s \leq f$ and $0 < c < 1$. Define:
\[ E_n = \{x \mid f_n(x) \geq c \cdot s(x)\} \]
Then $E_n \nearrow X$ since $f_n \nearrow f \geq s > cs$. Thus:
\[ \int f_n \geq \int_{E_n} f_n \geq c \int_{E_n} s \]
By continuity from below, $\int_{E_n} s \to \int_X s$. So $\lim \int f_n \geq c \int s$.

Taking $c \to 1$ and supremum over $s$: $\lim \int f_n \geq \int f$.
\end{proof}

\begin{corollary}
For measurable $f_n \geq 0$: $\int \sum_{n=1}^\infty f_n = \sum_{n=1}^\infty \int f_n$.
\end{corollary}

\begin{proof}
Let $S_N = \sum_{n=1}^N f_n$. Then $S_N \nearrow \sum_{n=1}^\infty f_n$. By MCT: $\int \sum f_n = \lim_N \int S_N = \lim_N \sum_{n=1}^N \int f_n = \sum_{n=1}^\infty \int f_n$.
\end{proof}

\section{Fatou's Lemma}

\begin{theorem}[Fatou's Lemma]
If $(f_n)$ are measurable with $f_n \geq 0$:
\[ \int \liminf_{n \to \infty} f_n \, d\mu \leq \liminf_{n \to \infty} \int f_n \, d\mu \]
\end{theorem}

\begin{proof}
Let $g_k = \inf_{n \geq k} f_n$. Then $g_k \nearrow \liminf f_n$ and $g_k \leq f_k$. By MCT:
\[ \int \liminf f_n = \lim_{k \to \infty} \int g_k \leq \liminf_{k \to \infty} \int f_k \]
\end{proof}

\section{Integrable Functions}

\begin{definition}
For measurable $f: X \to \bb{R}$, define $f^+ = \max(f, 0)$ and $f^- = \max(-f, 0)$.

$f$ is \textbf{integrable} ($f \in L^1(\mu)$) if $\int f^+ < \infty$ and $\int f^- < \infty$.

The integral is: $\int f \, d\mu = \int f^+ d\mu - \int f^- d\mu$.
\end{definition}

\begin{proposition}
$f \in L^1(\mu)$ iff $\int |f| d\mu < \infty$. Moreover: $\left|\int f\right| \leq \int |f|$.
\end{proposition}

\begin{proof}
$|f| = f^+ + f^-$, so $\int |f| = \int f^+ + \int f^- < \infty$ iff both are finite, iff $f \in L^1$.

Triangle inequality: $|\int f| = |\int f^+ - \int f^-| \leq \int f^+ + \int f^- = \int |f|$.
\end{proof}

\section{Dominated Convergence Theorem}

\begin{theorem}[DCT / Lebesgue]
Let $(f_n)$ be measurable with $f_n \to f$ \almosteverywhere{} If there exists $g \in L^1(\mu)$ with $|f_n| \leq g$ \almosteverywhere{} for all $n$, then:
\[ f \in L^1(\mu) \quad \text{and} \quad \lim_{n \to \infty} \int |f_n - f| d\mu = 0 \]
In particular: $\lim \int f_n = \int f$.
\end{theorem}

\begin{proof}
Since $|f| = \lim |f_n| \leq g$ \almosteverywhere{}, we have $f \in L^1$.

Apply Fatou to $g - f_n \geq 0$ and $g + f_n \geq 0$:
\[ \int g - f \leq \liminf \int (g - f_n) = \int g - \limsup \int f_n \]
\[ \int g + f \leq \liminf \int (g + f_n) = \int g + \liminf \int f_n \]
Thus $\limsup \int f_n \leq \int f \leq \liminf \int f_n$, so $\lim \int f_n = \int f$.

For $\int |f_n - f| \to 0$: apply Fatou to $2g - |f_n - f| \geq 0$.
\end{proof}

\section{Riemann vs Lebesgue}

\begin{theorem}[Lebesgue Criterion]
A bounded function $f: [a, b] \to \bb{R}$ is Riemann integrable iff $f$ is continuous \almosteverywhere{} (the set of discontinuities has Lebesgue measure 0).

If $f$ is Riemann integrable, then $f$ is Lebesgue integrable and:
\[ \int_a^b f(x) dx = \int_{[a,b]} f \, d\lambda \]
\end{theorem}

\begin{proof}
Let $D_f = \{x \in [a,b] \mid f \text{ is discontinuous at } x\}$.

\textbf{Oscillation:} For $x \in [a,b]$, define the oscillation:
\[ \omega_f(x) = \lim_{\delta \to 0^+} \left( \sup_{|y-x|<\delta} f(y) - \inf_{|y-x|<\delta} f(y) \right) \]

Then $x \in D_f$ iff $\omega_f(x) > 0$. Note $D_f = \bigcup_{n=1}^\infty \{x \mid \omega_f(x) \geq 1/n\}$.

\textbf{($\Rightarrow$):} Assume $f$ is Riemann integrable. For any $\epsilon > 0$, there exists partition $P$ with $\overline{S}(f,P) - \underline{S}(f,P) < \epsilon$.

Let $A_\epsilon = \{x \mid \omega_f(x) \geq \epsilon\}$. The contribution to $\overline{S} - \underline{S}$ from intervals meeting $A_\epsilon$ is $\geq \epsilon \cdot \lambda^*(A_\epsilon \setminus \text{endpoints})$.

Thus $\lambda^*(A_\epsilon) \leq (\overline{S} - \underline{S})/\epsilon \to 0$ as partition refines. So $\lambda(A_\epsilon) = 0$ for all $\epsilon > 0$, hence $\lambda(D_f) = 0$.

\textbf{($\Leftarrow$):} Assume $\lambda(D_f) = 0$. Given $\epsilon > 0$, cover $D_f$ by intervals of total length $< \epsilon$. On the compact set outside this cover, $f$ is continuous hence uniformly continuous.

Choose partition fine enough that oscillation on each subinterval is $< \epsilon$ (except those meeting the cover). Then $\overline{S} - \underline{S} < \epsilon(b-a) + 2M\epsilon$ where $M = \sup|f|$.

\textbf{Equality of integrals:} For Riemann integrable $f$, the upper/lower Riemann sums are Lebesgue integrals of simple functions that converge \almosteverywhere{} to $f$. By DCT (with dominating function $M$), the Lebesgue integral equals the Riemann integral.
\end{proof}


\chapter{Product Measures}

\section{Product $\sigma$-Algebras}

\begin{definition}
Let $(X, \mathcal{M})$ and $(Y, \mathcal{N})$ be measurable spaces. The \textbf{product $\sigma$-algebra} is:
\[ \mathcal{M} \otimes \mathcal{N} = \GenSigmaAlg{\{A \times B \mid A \in \mathcal{M}, B \in \mathcal{N}\}} \]
\end{definition}

\begin{proposition}
$\BorelAlg{\bb{R}^{n+m}} = \BorelAlg{\bb{R}^n} \otimes \BorelAlg{\bb{R}^m}$.
\end{proposition}

\begin{proof}
Let $\mathcal{R} = \{A \times B \mid A \in \BorelAlg{\bb{R}^n}, B \in \BorelAlg{\bb{R}^m}\}$. We show $\GenSigmaAlg{\mathcal{R}} = \BorelAlg{\bb{R}^{n+m}}$.

($\subseteq$): The projections $\pi_1: \bb{R}^{n+m} \to \bb{R}^n$ and $\pi_2: \bb{R}^{n+m} \to \bb{R}^m$ are continuous, hence measurable. For $A \in \BorelAlg{\bb{R}^n}$, $B \in \BorelAlg{\bb{R}^m}$: $A \times B = \pi_1^{-1}(A) \cap \pi_2^{-1}(B) \in \BorelAlg{\bb{R}^{n+m}}$.

($\supseteq$): Open sets in $\bb{R}^{n+m}$ are countable unions of open rectangles. Each open rectangle is in $\GenSigmaAlg{\mathcal{R}}$ since open sets are Borel.
\end{proof}

\section{Product Measures}

\begin{theorem}
Let $(X, \mathcal{M}, \mu)$ and $(Y, \mathcal{N}, \nu)$ be $\sigma$-finite measure spaces. There exists a unique measure $\mu \otimes \nu$ on $\mathcal{M} \otimes \mathcal{N}$ such that:
\[ (\mu \otimes \nu)(A \times B) = \mu(A) \cdot \nu(B) \]
\end{theorem}

\begin{proof}
\textbf{Existence:} For $E \in \mathcal{M} \otimes \mathcal{N}$, define:
\[ (\mu \otimes \nu)(E) = \int_X \nu(E_x) d\mu(x) \]
where $E_x = \{y \in Y \mid (x, y) \in E\}$ is the $x$-section. One shows $x \mapsto \nu(E_x)$ is measurable and this defines a measure.

\textbf{Uniqueness:} The rectangles $\{A \times B\}$ form a $\pi$-system generating $\mathcal{M} \otimes \mathcal{N}$. Apply the uniqueness theorem for $\sigma$-finite measures.
\end{proof}

\section{Tonelli's Theorem}

\begin{theorem}[Tonelli]
Let $f: X \times Y \to [0, \infty]$ be $(\mathcal{M} \otimes \mathcal{N})$-measurable. Then:
\begin{enumerate}
    \item For each $x$, $y \mapsto f(x, y)$ is $\mathcal{N}$-measurable
    \item $x \mapsto \int_Y f(x, y) d\nu(y)$ is $\mathcal{M}$-measurable
    \item The iterated integrals equal the product integral:
    \[ \int_{X \times Y} f \, d(\mu \otimes \nu) = \int_X \left(\int_Y f(x, y) d\nu(y)\right) d\mu(x) = \int_Y \left(\int_X f(x, y) d\mu(x)\right) d\nu(y) \]
\end{enumerate}
\end{theorem}

\begin{proof}
\textbf{Step 1: Characteristic functions of rectangles.}

Let $f = \CharFunc{A \times B}$ where $A \in \mathcal{M}$, $B \in \mathcal{N}$. Then:
\[ \int_Y f(x, y) d\nu(y) = \int_Y \CharFunc{A}(x) \CharFunc{B}(y) d\nu(y) = \CharFunc{A}(x) \nu(B) \]
So $x \mapsto \int_Y f d\nu$ is measurable. The iterated integral is:
\[ \int_X \CharFunc{A}(x) \nu(B) d\mu(x) = \mu(A) \nu(B) = (\mu \otimes \nu)(A \times B) = \int f d(\mu \otimes \nu) \]

\textbf{Step 2: Simple functions.}

By linearity, the result holds for simple functions $s = \sum_{i=1}^n c_i \CharFunc{E_i}$ where $E_i \in \mathcal{M} \otimes \mathcal{N}$. (Each $E_i$ can be approximated by finite disjoint unions of rectangles.)

\textbf{Step 3: General positive measurable functions.}

Let $f \geq 0$ be measurable. Choose simple $s_n \nearrow f$. By Step 2:
\[ \int s_n d(\mu \otimes \nu) = \int_X \left(\int_Y s_n(x,y) d\nu\right) d\mu \]

By MCT applied twice (inner then outer integral):
\[ \int f d(\mu \otimes \nu) = \lim_n \int s_n d(\mu \otimes \nu) = \lim_n \int_X \left(\int_Y s_n d\nu\right) d\mu = \int_X \left(\int_Y f d\nu\right) d\mu \]

Symmetry gives the other iterated integral.
\end{proof}

\section{Fubini's Theorem}

\begin{theorem}[Fubini]
Let $f \in L^1(\mu \otimes \nu)$. Then:
\begin{enumerate}
    \item For $\mu$-a.e. $x$: $y \mapsto f(x, y) \in L^1(\nu)$
    \item $x \mapsto \int_Y f(x, y) d\nu(y) \in L^1(\mu)$
    \item \[ \int_{X \times Y} f \, d(\mu \otimes \nu) = \int_X \left(\int_Y f(x, y) d\nu(y)\right) d\mu(x) \]
\end{enumerate}
\end{theorem}

\begin{proof}
\textbf{Step 1:} Apply Tonelli to $|f|$. Since $f \in L^1(\mu \otimes \nu)$:
\[ \int_{X \times Y} |f| d(\mu \otimes \nu) = \int_X \left(\int_Y |f(x,y)| d\nu\right) d\mu < \infty \]

\textbf{Step 2:} For $\mu$-a.e. $x$, the inner integral $\int_Y |f(x,y)| d\nu < \infty$.

Indeed, if $g(x) = \int_Y |f(x,y)| d\nu$, then $\int_X g d\mu < \infty$ implies $g < \infty$ $\mu$-a.e.

Thus $f_x = f(x, \cdot) \in L^1(\nu)$ for $\mu$-a.e. $x$.

\textbf{Step 3:} Decompose $f = f^+ - f^-$ and apply Tonelli to each:
\[ \int_{X \times Y} f^+ d(\mu \otimes \nu) = \int_X \left(\int_Y f^+(x,y) d\nu\right) d\mu \]
\[ \int_{X \times Y} f^- d(\mu \otimes \nu) = \int_X \left(\int_Y f^-(x,y) d\nu\right) d\mu \]

Both right-hand sides are finite (since $\leq \int |f| < \infty$), so we can subtract:
\[ \int f d(\mu \otimes \nu) = \int_X \left(\int_Y f d\nu\right) d\mu \]
\end{proof}


\chapter{$\Lp{p}$ Spaces}

\section{Definition}

\begin{definition}
For $1 \leq p < \infty$:
\[ \LpFunc{p}(\mu) = \left\{ f: X \to \bb{R} \text{ measurable} \mid \int |f|^p d\mu < \infty \right\} \]
The $\Lp{p}$ norm: $\LpNorm{f}{p} = \left(\int |f|^p d\mu\right)^{1/p}$.

$\Lp{p}(\mu) = \LpFunc{p}(\mu) / \sim$ where $f \sim g$ iff $f = g$ \almosteverywhere{}
\end{definition}

\section{Hölder's Inequality}

\begin{theorem}[Hölder]
Let $1 < p, q < \infty$ with $\frac{1}{p} + \frac{1}{q} = 1$. If $f \in \Lp{p}$ and $g \in \Lp{q}$:
\[ \LpNorm{fg}{1} \leq \LpNorm{f}{p} \LpNorm{g}{q} \]
\end{theorem}

\begin{proof}
WLOG $\LpNorm{f}{p} = \LpNorm{g}{q} = 1$. Use Young's inequality: $ab \leq \frac{a^p}{p} + \frac{b^q}{q}$.

For $a = |f(x)|$, $b = |g(x)|$:
\[ |f(x)g(x)| \leq \frac{|f(x)|^p}{p} + \frac{|g(x)|^q}{q} \]
Integrate: $\int |fg| \leq \frac{1}{p} + \frac{1}{q} = 1 = \LpNorm{f}{p} \LpNorm{g}{q}$.
\end{proof}

\section{Minkowski's Inequality}

\begin{theorem}[Minkowski]
For $1 \leq p < \infty$ and $f, g \in \Lp{p}$:
\[ \LpNorm{f + g}{p} \leq \LpNorm{f}{p} + \LpNorm{g}{p} \]
\end{theorem}

\begin{proof}
For $p = 1$: triangle inequality. For $p > 1$:
\[ |f + g|^p \leq |f + g|^{p-1}|f| + |f + g|^{p-1}|g| \]
Apply Hölder with exponents $p$ and $q = p/(p-1)$:
\[ \int |f + g|^p \leq \LpNorm{f + g}{p}^{p-1}(\LpNorm{f}{p} + \LpNorm{g}{p}) \]
Divide by $\LpNorm{f + g}{p}^{p-1}$.
\end{proof}

\section{Completeness}

\begin{theorem}[Riesz-Fischer]
$(\Lp{p}(\mu), \LpNorm{\cdot}{p})$ is a Banach space (complete normed vector space).
\end{theorem}

\begin{proof}
Let $(f_n)$ be Cauchy in $\Lp{p}$. Choose subsequence $(f_{n_k})$ with $\LpNorm{f_{n_{k+1}} - f_{n_k}}{p} < 2^{-k}$.

Define $g = \sum_{k=1}^\infty |f_{n_{k+1}} - f_{n_k}|$. By Minkowski: $\LpNorm{g}{p} \leq \sum 2^{-k} < \infty$.

So $g < \infty$ \almosteverywhere{}, meaning $\sum (f_{n_{k+1}} - f_{n_k})$ converges absolutely \almosteverywhere{}

Define $f = \lim f_{n_k}$ (exists \almosteverywhere{}). By Fatou: $f \in \Lp{p}$.

By DCT: $\LpNorm{f_{n_k} - f}{p} \to 0$. Since $(f_n)$ is Cauchy, $\LpNorm{f_n - f}{p} \to 0$.
\end{proof}


\chapter{Convergence in Measure}

\section{Definition}

\begin{definition}
$f_n \to f$ \textbf{in measure} if for every $\alpha > 0$:
\[ \lim_{n \to \infty} \mu(\{x \mid |f_n(x) - f(x)| \geq \alpha\}) = 0 \]
\end{definition}

\section{Relations Between Convergence Modes}

\begin{proposition}
\begin{enumerate}
    \item $\Lp{p}$ convergence $\implies$ convergence in measure
    \item A.e. convergence + finite measure $\implies$ convergence in measure
\end{enumerate}
\end{proposition}

\begin{proof}
(1) Chebyshev: $\mu(|f_n - f| \geq \alpha) \leq \frac{1}{\alpha^p}\|f_n - f\|_p^p \to 0$.

(2) By Egorov's theorem (uniform convergence on large sets).
\end{proof}

\section{Riesz's Theorem}

\begin{theorem}
If $f_n \to f$ in measure, there exists a subsequence $f_{n_k} \to f$ \almosteverywhere{}
\end{theorem}

\begin{proof}
Choose $n_k$ such that $\mu(|f_{n_k} - f| \geq 2^{-k}) < 2^{-k}$.

Let $E_k = \{|f_{n_k} - f| \geq 2^{-k}\}$. Then $\sum \mu(E_k) < \infty$.

By Borel-Cantelli: \almosteverywhere{} $x$ is in only finitely many $E_k$, so $f_{n_k}(x) \to f(x)$.
\end{proof}

\section{Vitali Convergence}

\begin{definition}
$(f_n)$ is \textbf{uniformly integrable} if for every $\epsilon > 0$, there exists $\delta > 0$ such that:
\[ \mu(A) < \delta \implies \sup_n \int_A |f_n| < \epsilon \]
\end{definition}

\begin{theorem}[Vitali]
Let $\mu(X) < \infty$ and $f_n \to f$ in measure. Then $f_n \to f$ in $L^1$ iff $(f_n)$ is uniformly integrable.
\end{theorem}

\begin{proof}
\textbf{($\Rightarrow$)} Suppose $f_n \to f$ in $L^1$. Given $\epsilon > 0$, for large $n$, $\|f_n - f\|_1 < \epsilon/3$.

Since $f \in L^1$, there exists $\delta > 0$ such that $\mu(A) < \delta \implies \int_A |f| < \epsilon/3$.

For the finitely many remaining $f_1, \ldots, f_N$, reduce $\delta$ so $\int_A |f_k| < \epsilon$ for $k \leq N$ when $\mu(A) < \delta$.

For $n > N$: $\int_A |f_n| \leq \int_A |f_n - f| + \int_A |f| < \epsilon/3 + \epsilon/3 < \epsilon$.

\textbf{($\Leftarrow$)} Suppose $(f_n)$ is uniformly integrable and $f_n \to f$ in measure.

By Riesz, some subsequence $f_{n_k} \to f$ \almosteverywhere{} By Fatou, $f \in L^1$.

Given $\epsilon > 0$, choose $\delta$ from uniform integrability. By Egorov, there exists $E$ with $\mu(E^c) < \delta$ and $f_n \to f$ uniformly on $E$.

Then $\int |f_n - f| = \int_E |f_n - f| + \int_{E^c} |f_n - f| \leq \mu(E) \cdot \sup_E |f_n - f| + \int_{E^c} |f_n| + \int_{E^c} |f|$.

The first term $\to 0$. The last two are $< 2\epsilon$ by uniform integrability.
\end{proof}


\chapter{Applications}

\section{Differentiation Under the Integral}

\begin{theorem}[Leibniz Rule]
Let $f: I \times X \to \bb{R}$ where $I \subseteq \bb{R}$ is an interval. If:
\begin{enumerate}
    \item $t \mapsto f(t, x)$ is differentiable for each $x$
    \item $|\frac{\partial f}{\partial t}(t, x)| \leq g(x)$ for some $g \in L^1(\mu)$
\end{enumerate}
Then $\phi(t) = \int_X f(t, x) d\mu(x)$ is differentiable and:
\[ \phi'(t) = \int_X \frac{\partial f}{\partial t}(t, x) d\mu(x) \]
\end{theorem}

\begin{proof}
Apply DCT to the difference quotients $\frac{f(t+h, x) - f(t, x)}{h}$.
\end{proof}

\section{Gaussian Integral}

\begin{theorem}
$\int_{\bb{R}} e^{-x^2} dx = \sqrt{\pi}$.
\end{theorem}

\begin{proof}
Let $I = \int_\bb{R} e^{-x^2} dx$. Then:
\[ I^2 = \int_{\bb{R}^2} e^{-(x^2+y^2)} dx\,dy = \int_0^{2\pi} \int_0^\infty e^{-r^2} r\,dr\,d\theta = 2\pi \cdot \frac{1}{2} = \pi \]
Thus $I = \sqrt{\pi}$.
\end{proof}


\chapter{Exercises}

\section{$\sigma$-Algebras and Borel Sets}

\begin{enumerate}
    \item Prove that $\BorelAlg{\bb{R}} = \GenSigmaAlg{\{(a,b) \mid a < b\}}$.
    \item Show $\bb{Q} \in \BorelAlg{\bb{R}}$ and $\bb{R} \setminus \bb{Q} \in \BorelAlg{\bb{R}}$.
    \item Let $\mathcal{F} = \{(a, \infty) \mid a \in \bb{R}\}$. Prove $\GenSigmaAlg{\mathcal{F}} = \BorelAlg{\bb{R}}$.
    \item Show that the collection of half-open intervals $\{(a, b] \mid a \leq b\}$ forms a $\pi$-system.
    \item Prove the Cantor set is Borel measurable with $\lambda(C) = 0$.
    \item Show that $\{f^{-1}(B) \mid B \in \BorelAlg{\bb{R}}\}$ is a $\sigma$-algebra for any $f: X \to \bb{R}$.
    \item Prove that every $\pi$-system closed under complements is an algebra.
    \item Show that $\GenSigmaAlg{\mathcal{F}_1 \cup \mathcal{F}_2} = \GenSigmaAlg{\GenSigmaAlg{\mathcal{F}_1} \cup \mathcal{F}_2}$.
    \item Prove that $|\BorelAlg{\bb{R}}| = |\bb{R}|$ (cardinality of continuum).
    \item Show that $\BorelAlg{\bb{R}^n} = \GenSigmaAlg{\text{open balls}}$.
\end{enumerate}

\section{Measures}

\begin{enumerate}
    \item Prove $\sigma$-subadditivity from $\sigma$-additivity using disjointification.
    \item Let $A_n = [n, \infty)$. Show $\lim \lambda(A_n) \neq \lambda(\cap A_n)$. Why doesn't continuity from above apply?
    \item Prove continuity from below directly from $\sigma$-additivity.
    \item Let $\mu$ be counting measure on $\bb{N}$. Find $\mu(\{n \mid n \text{ is even}\})$.
    \item Show that if $\mu(A_n) = 0$ for all $n$, then $\mu(\cup A_n) = 0$.
    \item Prove that $\mu(A \triangle B) = 0$ defines an equivalence relation.
    \item Show that $\mu(\liminf A_n) \leq \liminf \mu(A_n)$ (Fatou for sets).
    \item Prove Borel-Cantelli: if $\sum \mu(A_n) < \infty$, then $\mu(\limsup A_n) = 0$.
    \item Let $\mu$ be finite. Prove $\mu(\limsup A_n) \geq \limsup \mu(A_n)$ is false in general.
    \item Construct a measure on $\bb{N}$ that is not $\sigma$-finite.
    \item Prove that Dirac measure $\delta_x$ is a measure.
    \item Show that $\mu(A) = \sum_{n \in A} 2^{-n}$ is a finite measure on $\bb{N}$.
\end{enumerate}

\section{Carathéodory and Lebesgue Measure}

\begin{enumerate}
    \item Verify that every interval $[a, b]$ is $\lambda^*$-measurable.
    \item Prove $\lambda^*(A) = \inf\{\lambda(G) \mid G \text{ open}, G \supseteq A\}$.
    \item Construct a Lebesgue measurable set that is not Borel.
    \item Show $\BorelAlg{\bb{R}} \subsetneq \text{Leb}(\bb{R}) \subsetneq \mathcal{P}(\bb{R})$.
    \item Prove that the completion of Borel measure is Lebesgue measure.
    \item Show $\lambda([0,1] \cap \bb{Q}) = 0$ and $\lambda([0,1] \setminus \bb{Q}) = 1$.
    \item Prove that $\lambda$ is translation-invariant: $\lambda(A + x) = \lambda(A)$.
    \item Show $\lambda(cA) = |c|^n \lambda(A)$ for $A \subseteq \bb{R}^n$.
    \item Construct a non-measurable set using the axiom of choice.
    \item Prove that every set of positive measure contains a non-measurable subset.
    \item Show that $\lambda^*(A) = 0$ iff $A$ is contained in a Borel null set.
    \item Prove Lebesgue's density theorem: $\lim_{r \to 0} \frac{\lambda(A \cap B_r(x))}{\lambda(B_r(x))} = 1$ for \almosteverywhere{} $x \in A$.
\end{enumerate}

\section{Riemann vs Lebesgue}

\begin{enumerate}
    \item Prove the Dirichlet function $\CharFunc{\bb{Q} \cap [0,1]}$ is Lebesgue integrable with $\int = 0$.
    \item Show it is not Riemann integrable.
    \item Find a function continuous \almosteverywhere{} that is Riemann integrable.
    \item Construct $f$ bounded with uncountable discontinuity set but still Riemann integrable.
    \item Prove: Riemann integrable $\implies$ Lebesgue integrable with same value.
    \item Show that $f(x) = \sin(1/x)$ on $(0,1]$, $f(0) = 0$ is Riemann integrable.
    \item Prove the Thomae function is Riemann integrable.
    \item Show that $f_n(x) = x^n$ on $[0,1]$ converges pointwise but $\int_0^1 f_n \not\to \int_0^1 \lim f_n$ in Riemann sense.
    \item Compute $\int_0^1 \CharFunc{C} dx$ where $C$ is the Cantor set.
    \item Show that if $f$ is monotone on $[a,b]$, then $f$ is Riemann integrable.
\end{enumerate}

\section{Measurable Functions}

\begin{enumerate}
    \item Show that continuous functions are Borel measurable.
    \item Prove that monotone functions on $\bb{R}$ are Borel measurable.
    \item Show $f \cdot g$ measurable when $f, g$ measurable (direct proof).
    \item Prove that $\{f = g\}$ is measurable when $f, g$ are measurable.
    \item Show that $\limsup f_n$ is measurable when each $f_n$ is measurable.
    \item Prove that the composition of measurable functions is measurable.
    \item Show that $f^{-1}(\{y\})$ is measurable for each $y$ iff $f$ is measurable (false in general).
    \item Prove Lusin's theorem: measurable $f$ is continuous on a large set.
    \item Show that every measurable $f \geq 0$ is a limit of simple functions.
    \item Prove Egorov's theorem: \almosteverywhere{} convergence implies almost uniform convergence on finite measure spaces.
\end{enumerate}

\section{Convergence Theorems}

\begin{enumerate}
    \item Compute $\lim_{n \to \infty} \int_0^1 nx^{n-1} dx$. Does MCT apply? DCT?
    \item Let $f_n = n \CharFunc{[0, 1/n]}$. Show $f_n \to 0$ \almosteverywhere{} but $\int f_n = 1$.
    \item Use MCT to prove $\int \sum f_n = \sum \int f_n$ for $f_n \geq 0$.
    \item Apply Fatou to find $\liminf \int f_n$ when $f_n = \CharFunc{[n, n+1]}$.
    \item Use DCT to compute $\lim \int_0^\infty \frac{\sin(x/n)}{x} e^{-x} dx$.
    \item Compute $\lim_{n \to \infty} \int_0^n (1 - x/n)^n e^{x/2} dx$.
    \item Show $\lim_{n \to \infty} \int_0^\infty \frac{n \sin(x/n)}{x(1+x^2)} dx = \pi/2$.
    \item Prove: if $f_n \to f$ in $L^1$ and $f_n \to g$ \almosteverywhere{}, then $f = g$ \almosteverywhere{}
    \item Show that MCT fails without monotonicity.
    \item Compute $\lim_{n \to \infty} \int_0^1 \frac{nx}{1+n^2x^2} dx$.
    \item Prove the reverse Fatou lemma with additional hypothesis.
    \item Show $\int_0^\infty e^{-x^2} dx = \sqrt{\pi}/2$ using Fubini.
\end{enumerate}

\section{$\Lp{p}$ Spaces}

\begin{enumerate}
    \item Let $f_n = n^{1/p} \CharFunc{[1/n, 2/n]}$. Show $\|f_n\|_p = 1$ but $f_n \to 0$ pointwise.
    \item Prove Young's inequality: $ab \leq \frac{a^p}{p} + \frac{b^q}{q}$.
    \item Use Hölder to show $\|f\|_1 \leq \|f\|_2 \mu(X)^{1/2}$ for finite measure.
    \item Prove $L^2$ convergence implies $L^1$ convergence on finite measure spaces.
    \item Find $f \in L^1([0,1])$ but $f \notin L^2([0,1])$.
    \item Show $\\Lp{p} \subseteq \Lp{q}$ for $p > q$ on finite measure spaces.
    \item Prove that $L^2$ is a Hilbert space with inner product $\langle f, g \rangle = \int f \bar{g}$.
    \item Show that simple functions are dense in $\Lp{p}$.
    \item Prove that $C_c(\bb{R})$ is dense in $\Lp{p}(\bb{R})$ for $p < \infty$.
    \item Find the dual exponent $q$ for $p = 3/2$.
    \item Show $\|f\|_\infty = \lim_{p \to \infty} \|f\|_p$ on finite measure spaces.
    \item Prove that $L^\infty$ is not separable on $[0,1]$.
\end{enumerate}

\section{Product Measures}

\begin{enumerate}
    \item Use Fubini to compute $\int_0^1 \int_0^1 \frac{x-y}{(x+y)^3} dy\, dx$.
    \item Compute $\int_{\bb{R}^2} e^{-(x^2+y^2)} dx\, dy$ using polar coordinates.
    \item Show iterated integrals may differ if $f \notin L^1$.
    \item Prove $\BorelAlg{\bb{R}^2} = \BorelAlg{\bb{R}} \otimes \BorelAlg{\bb{R}}$.
    \item Compute $\int_0^\infty \frac{\sin x}{x} dx$ using parameter integrals.
    \item Use Fubini to prove $\int_0^\infty \frac{e^{-ax} - e^{-bx}}{x} dx = \ln(b/a)$.
    \item Compute $\int_0^1 \int_0^1 \frac{1}{1-xy} dx\, dy$ (hint: expand as series).
    \item Show that $(\mu \otimes \nu) \otimes \rho = \mu \otimes (\nu \otimes \rho)$.
    \item Prove Cavalieri's principle for Lebesgue measure.
    \item Compute the volume of the $n$-ball using Fubini.
\end{enumerate}

\section{Convergence in Measure}

\begin{enumerate}
    \item Show $\Lp{p}$ convergence implies convergence in measure.
    \item Find $(f_n)$ converging in measure but not \almosteverywhere{}
    \item Prove Riesz: convergence in measure has \almosteverywhere{} convergent subsequence.
    \item Define uniform integrability. Give an example.
    \item State and apply Vitali's theorem.
    \item Show that \almosteverywhere{} convergence does not imply convergence in measure (infinite measure).
    \item Prove: $f_n \to f$ in measure and $g_n \to g$ in measure implies $f_n + g_n \to f + g$ in measure.
    \item Show that $f_n \to f$ in measure iff every subsequence has a further subsequence converging \almosteverywhere{}
    \item Prove that uniform integrability + convergence in measure implies $L^1$ convergence.
    \item Give an example where $f_n \to f$ in $L^1$ but not uniformly.
\end{enumerate}

\section{Parameter Integrals}

\begin{enumerate}
    \item Prove $\phi(t) = \int_0^\infty e^{-tx} \sin x \, dx$ is differentiable for $t > 0$.
    \item Compute $\int_0^\infty \frac{e^{-x} - e^{-2x}}{x} dx$ using parameter integrals.
    \item Prove $\Gamma(n+1) = n!$ where $\Gamma(s) = \int_0^\infty x^{s-1} e^{-x} dx$.
    \item Justify differentiation under the integral for $\int_0^1 x^t dx$.
    \item Use Leibniz rule to compute $\frac{d}{dt} \int_0^t f(t, x) dx$.
    \item Compute $\int_0^\infty \frac{\ln x}{x^2 - 1} dx$.
    \item Prove $\Gamma(1/2) = \sqrt{\pi}$.
    \item Use parameter integrals to compute $\int_0^\infty \frac{\sin x}{x} dx = \pi/2$.
    \item Prove the beta function identity: $B(a,b) = \frac{\Gamma(a)\Gamma(b)}{\Gamma(a+b)}$.
    \item Compute $\int_0^\infty x^n e^{-x^2} dx$ in terms of $\Gamma$.
\end{enumerate}

\section{Modes of Convergence}

\begin{enumerate}
    \item Construct $(f_n)$ on $[0,1]$ with $f_n \to 0$ in measure but $f_n(x) \not\to 0$ for any $x$.
    \item Show: \almosteverywhere{} convergence + domination implies $\Lp{p}$ convergence.
    \item Prove that on finite measure spaces, $\Lp{p}$ convergence implies $\Lp{q}$ convergence for $q < p$.
    \item Show that uniform convergence implies $\Lp{p}$ convergence on finite measure spaces.
    \item Prove that $f_n \to f$ uniformly iff $\|f_n - f\|_\infty \to 0$.
    \item Show the "typewriter sequence" converges in measure but not \almosteverywhere{}
    \item Prove: $L^1$ convergence + \almosteverywhere{} convergence to $g$ implies $f = g$ \almosteverywhere{}
    \item Give an example of $\Lp{p}$ convergence without \almosteverywhere{} convergence.
\end{enumerate}

\section{Simple Functions and Integration}

\begin{enumerate}
    \item Show that $\int (f + g) = \int f + \int g$ for integrable $f, g$.
    \item Prove that $\int cf = c \int f$ for $c \in \bb{R}$.
    \item Show that $f \leq g$ \almosteverywhere{} implies $\int f \leq \int g$.
    \item Prove that $|\int f| \leq \int |f|$ (triangle inequality).
    \item Show that $\int f = 0$ for all measurable $E$ implies $f = 0$ \almosteverywhere{}
    \item Prove Chebyshev's inequality: $\mu(\{|f| \geq \alpha\}) \leq \frac{1}{\alpha} \int |f|$.
    \item Show that if $\int |f|^p < \infty$ and $\int |f|^q < \infty$, then $\int |f|^r < \infty$ for $p < r < q$.
    \item Prove Markov's inequality: $\mu(\{f \geq \alpha\}) \leq \frac{\int f}{\alpha}$ for $f \geq 0$.
\end{enumerate}

